\section{Does binding depend on one epitope island or both?}\label{sec:islands}
A two-island epitope binds its antibody through two spatially separate patches, and the native
structure alone cannot say how the interaction is divided between them. Does the antibody need
island~A, island~B, or only both together? The crystal complex shows the bound state but not
which contacts are load-bearing, and no purely computational metric on the native geometry can
answer it. The way to find out is to build the reagents that separate the two contributions:
scaffold island~A on its own, island~B on its own, and the two in their native arrangement, then
let the PepSeq assay read off which presentations the antibody still binds. If a single island
suffices, it carries the interaction; if neither binds alone, the epitope is genuinely
discontinuous and both islands are required. Different two-island epitopes may distribute their
binding differently, and that distribution is exactly what the per-island run
(Section~\ref{sec:perisland}) is built to measure.

\paragraph{A related, computational question: are two islands harder to scaffold?}
Before turning to binding, the design data answer a narrower question on their own: holding the
native geometry fixed, are two-island epitopes harder to scaffold than one? This is the
island-\emph{count} control, a first cut at the within-DP3 stratification of
Section~\ref{sec:open}. We split the DP3 design set (RFD1+MPNN, \texttt{dp2.parquet}) by
epitope island count, where an island is a spatially contiguous patch of epitope residues and
the count is the number of fixed segments in the contig (the \texttt{epitope\_chunks} field).
Of the 59 epitopes, 13 are single-island and 46 are two-island; none have more under this
grouping. Restricting to designs with an AF3 result, that is 549 single-island and
\num{123813} two-island designs.

Across every four-filter metric the single-island designs sit at better values
(Fig.~\ref{fig:islands}): median epitope-chunk RMSD 2.22 vs \SI{3.86}{\angstrom}, median
global PAE 9.5 vs 11.9, and fewer AF3 clashes. The binary pass rate runs the other way ---
single-island $0.24\%$ (\num{2}/\num{832}) vs two-island $0.56\%$ (\num{838}/\num{150400}), on
the same per-design denominator as Section~\ref{sec:rfd} --- but that is misleading: of the
\num{838} two-island passes, \num{630} come from a single easy outlier (\texttt{7ox3\_0P}) and
another \num{149} from \texttt{5fhx\_1P}, so two epitopes account for \num{779} of them. Strip
the outlier and the two-island rate falls below single-island ($0.17\%$ vs $0.36\%$ per valid
prediction); the single-island set has no comparable outlier and is broadly easier. So on the
distributions, fewer islands is easier to scaffold, consistent with the multi-island
difficulty argument in Section~\ref{sec:intro}; the binary rate just hides it. What remains for
the clean control is to split further by island \emph{size}, since the single-island epitopes
are also shorter on average.

So the two questions point the same way for the design itself --- multi-island is both
biologically ambiguous and harder to scaffold cleanly --- and both motivate building each
island as a standalone, individually testable construct.

\begin{figure}[h]\centering
\figorbox{island_split_metrics.png}{1.0}
\caption{DP3 design metrics (RFD1+MPNN, \texttt{dp2.parquet}) split by epitope island count:
\textcolor[HTML]{1f77b4}{blue} = single-island ($n=549$),
\textcolor[HTML]{d62728}{red} = two-island ($n=\num{123813}$); only designs with an AF3
result are included. Panels, left to right: epitope-chunk RMSD, overall RMSD, global mean PAE,
and number of AF3 clashing residues, the four selection metrics. Distributions are
area-normalized (density) because the groups differ in size by $\sim$200$\times$. Dotted
vertical lines mark each group's median; dashed grey lines mark the four-filter thresholds
(epitope RMSD $\le 1$, overall RMSD $\le 2$, PAE $< 5$, clashes $= 0$). Single-island designs
are shifted toward better values on every metric (median epitope RMSD 2.22 vs
\SI{3.86}{\angstrom}, median PAE 9.5 vs 11.9). Regenerated by
\texttt{episcaf\_analysis/viz/plot\_island\_split.py}.}
\label{fig:islands}
\end{figure}
